\chapter{PENDAHULUAN}

\section{Latar Belakang}
% Latar belakang memuat uraian alasan mengapa Anda memilih topik ini, dilengkapi data relevan yang menunjang[cite: 308].
% Jika ada penelitian serupa sebelumnya, sebutkan perbedaannya dengan yang Anda buat atau sebutkan kelemahan penelitian terdahulu[cite: 309, 310].
Tuliskan latar belakang penelitian Anda di sini...

\section{Rumusan Masalah}
% Berisi pernyataan yang jelas, spesifik, dan terfokus mengenai inti masalah[cite: 311]. 
% Ditulis secara singkat, padat, dan jelas[cite: 312].
Berdasarkan latar belakang di atas, rumusan masalah dalam penelitian ini adalah...

\section{Ruang Lingkup}
% Memberikan batasan yang jelas dari rumusan masalah agar bahasan lebih tajam (misal: batasan data, metode, tools, dll)[cite: 313].
Ruang lingkup pada penelitian ini dibatasi pada...

\section{Tujuan Penelitian}
% Menggambarkan target capaian yang diharapkan, ditulis dalam naratif atau poin/paragraf, dan harus sesuai dengan Rumusan Masalah[cite: 314].
Tujuan dari penelitian ini adalah...

\section{Sistematika Penulisan}
% Memberikan gambaran umum dari bab ke bab secara narasi (jangan ditulis dalam format poin-poin/bullet)[cite: 315, 316].
Penulisan skripsi ini disusun dengan sistematika sebagai berikut. Bab 1 Pendahuluan membahas mengenai...